\chapter{Informed Search}\label{ch:informed}
\begin{center}\small\textit{Uninformed search explores blindly; informed search guesses wisely\\--- guided by a heuristic that estimates how far remains to the goal.}\end{center}

\noindent\fbox{\parbox{\dimexpr\linewidth-2\fboxsep-2\fboxrule}{%
\textbf{Learning outcomes.} After this chapter you will be able to: (1)~define $g$, $h$, $f$ and explain $f=g+h$; (2)~distinguish greedy best-first from A$^{*}$; (3)~test admissibility vs.\ consistency; (4)~prove A$^{*}$ optimality; (5)~run IDA$^{*}$ and choose a heuristic by dominance.}}
\vspace{0.3em}

\noindent\textit{Prerequisites:} uninformed search (Ch.~2), priority queues, and big-$O$. We build from zero: no prior heuristic knowledge assumed.

% Notation: states are nodes, actions are edges with cost c(n,n')>0, g is path cost, h is heuristic, f=g+h.

\section{From Blind to Guided Search}
Uninformed methods (\textsc{BFS}, \textsc{DFS}, UCS) use only the path cost $g(n)$ --- the cheapest known cost from start $S$ to node $n$. They cannot tell whether $n$ is promising. \emph{Informed} (heuristic) search adds a guess $h(n)\approx h^{*}(n)$, the true cheapest cost $n\to$ goal. The art is choosing $h$ that is informative yet never misleading.

\begin{definition}[Heuristic function]
A heuristic $h:\mathcal{S}\to\mathbb{R}_{\ge 0}$ maps each state to a non-negative estimate of the cheapest cost to the nearest goal, with $h(G)=0$ for every goal $G$.
\end{definition}

Intuition: on a road map, \emph{straight-line distance} to Bucharest underestimates any real route --- it is optimistic and therefore safe. Manhattan distance on a 4-neighbour grid has the same property when each move costs $1$. Relaxing the problem (ignore roads, ignore obstacles) is the standard way to invent admissible heuristics.

\section{Greedy Best-First Search}
Greedy best-first expands the node that \emph{looks} closest to goal:
\[
  f(n)=h(n).
\]
It picks the smallest $h$, ignoring how far we have already travelled. It is fast and often finds a solution quickly, but it is neither complete (in infinite spaces) nor optimal.

\begin{example}[Greedy trap]
On the Romania map, straight-line $h$ makes greedy go Arad $\to$ Sibiu $\to$ Fagaras $\to$ Bucharest ($450$\,km) even though Arad $\to$ Sibiu $\to$ Rimnicu $\to$ Pitesti $\to$ Bucharest ($418$\,km) is cheaper. Greedy never looks back at $g$.
\end{example}

UCS is the opposite extreme: $f(n)=g(n)$, optimal but explores equally in all directions. \textsc{A}$^{*}$ combines both.

\section{A\texorpdfstring{$^{*}$}{*} Search: \texorpdfstring{$f(n)=g(n)+h(n)$}{f(n)=g(n)+h(n)}}
\begin{definition}[A$^{*}$ evaluation]
For node $n$, let $g(n)$ be the cheapest known cost $S\to n$ and $h(n)$ the heuristic. A$^{*}$ orders its frontier (priority queue) by
\[
  f(n)=g(n)+h(n),
\]
an estimate of the total cost of the best solution constrained to go through $n$.
\end{definition}
At each step A$^{*}$ pops the smallest-$f$ node, tests for goal, expands successors, and pushes them with their $f$ values. With a min-heap this is $O(\log |\text{frontier}|)$ per push/pop. If $h$ never overestimates, the first goal popped is optimal.

\begin{figure}[h]\centering
\begin{tikzpicture}[font=\scriptsize, every node/.style={font=\scriptsize},
  level distance=13mm, sibling distance=22mm,
  edge from parent/.style={draw, -{Stealth}, thick, black!60}]
\tikzset{st/.style={draw, rounded corners, align=center, minimum width=1.65cm, minimum height=7mm, fill=blue!10, draw=blue!50!black},
  expand/.style={fill=orange!22, draw=orange!60!black, thick},
  goal/.style={fill=green!20, draw=green!50!black, thick},
  fringe/.style={fill=violet!10, draw=violet!40!black, dashed}}
\node[st, fill=yellow!18] (root) {S\\$g\!=\!0,h\!=\!6,f\!=\!6$}
  child {node[st,expand] {A\\$g\!=\!1,h\!=\!4,f\!=\!5$}
    child {node[st,fringe] {C\\$1\!+\!3\!+\!2\!=\!6$}
      child {node[st,goal] {G\\$5\!+\!0\!=\!5^{\star}$} edge from parent node[midway,left,font=\tiny]{2}}
      edge from parent node[midway,left,font=\tiny]{3}}
    child {node[st,fringe] {D\\$1\!+\!4\!+\!1\!=\!6$} edge from parent node[midway,right,font=\tiny]{4}}
    edge from parent node[midway,left,font=\tiny]{1}}
  child {node[st] {B\\$g\!=\!4,h\!=\!2,f\!=\!6$}
    child {node[st,fringe] {G$_{2}$\\$4\!+\!5\!+\!0\!=\!9$} edge from parent node[midway,right,font=\tiny]{5}}
    edge from parent node[midway,right,font=\tiny]{4}};
\node[below=1mm of root, font=\tiny, text=gray!60!black] {frontier ordered by $f$};
\draw[->, thick, red!60!black] (3.2,-1.2) -- (2.1,-1.9) node[midway,above,sloped,font=\tiny,fill=white]{expand $f=5$ first};
\node[draw, fill=white, rounded corners, font=\tiny, align=left] at (5.8,-3.2) {\textbf{Legend:}\\{\color{orange!70!black}$\blacksquare$} expanded\\{\color{violet!60!black}$\blacksquare$} fringe\\{\color{green!50!black}$\blacksquare$} goal};
\end{tikzpicture}
\caption{A$^{*}$ search tree. Each node shows $g+h=f$. A$^{*}$ expands smallest $f$ first; the goal via A--C ($f=5$) is found before the costlier B route ($f=9$).}
\label{fig:astar-tree}
\end{figure}

\section{Admissibility and Consistency}
\begin{definition}[Admissibility]
$h$ is \emph{admissible} iff it never overestimates: $0\le h(n)\le h^{*}(n)$ for all $n$, where $h^{*}(n)$ is the true cheapest cost $n\to$ goal.
\end{definition}
\begin{definition}[Consistency / Monotonicity]
$h$ is \emph{consistent} iff for every edge $n\to n'$ with cost $c(n,n')$,
\[
  h(n)\le c(n,n')+h(n')\quad\text{(triangle inequality),}
\]
and $h(G)=0$. Consistency implies admissibility; the converse is false.
\end{definition}
Consistency guarantees $f$ is non-decreasing along any path, so A$^{*}$ with graph search never needs to reopen a closed node. With tree search, admissibility alone suffices for optimality; with graph search, consistency is needed (otherwise a better $g$ to a closed node may appear later).

\begin{figure}[h]\centering
\begin{tikzpicture}[font=\scriptsize]
\draw[->, thick, black!50] (-0.3,0) -- (7.2,0) node[right] {cost};
\fill[green!55!black] (0,0) circle (1.8pt) node[below=2pt, font=\tiny] {$0$};
\fill[blue!60!black] (5.0,0) circle (1.8pt) node[below=2pt, font=\tiny] {$h^{*}(n)$};
\fill[orange!80!black] (3.2,0) circle (1.8pt) node[below=2pt, font=\tiny] {$h(n)$};
\draw[thick, green!40!black] (0,0.55) -- (5.0,0.55);
\draw[thick, orange!70!black] (0,0.32) -- (3.2,0.32);
\node[font=\tiny, green!40!black] at (2.5,0.78) {true cost $h^{*}$};
\node[font=\tiny, orange!70!black] at (1.6,0.12) {admissible $h\le h^{*}$};
\draw[->, red!65!black, dashed, thick] (5.0,0.15) -- (6.4,0.15) node[midway,above,font=\tiny,fill=white]{overestimate $\Rightarrow$ inadmissible};
\fill[red!65!black] (6.4,0) circle (1.8pt) node[below=2pt,font=\tiny] {$h$};
\begin{scope}[yshift=-1.40cm]
\node[draw, fill=blue!10, circle, minimum size=7mm] (a) at (0,0) {$n$};
\node[draw, fill=blue!10, circle, minimum size=7mm] (b) at (3.0,0) {$n'$};
\node[draw, fill=green!15, circle, minimum size=7mm] (g) at (6.0,0) {$G$};
\draw[->, thick, blue!60!black] (a) -- (b) node[midway,above] {$c=2$};
\draw[->, thick, blue!60!black] (b) -- (g) node[midway,above] {$h(n')=3$};
\draw[->, thick, violet!60!black, bend left=22] (a) to node[above] {$h(n)\le 2+3=5$} (g);
\node[below=2mm of a, font=\tiny, text=gray!60!black] {$h(n)=4$};
\node[below=2mm of b, font=\tiny, text=gray!60!black] {$h(n')=3$};
\node[font=\tiny, fill=yellow!12, draw=orange!40!black, rounded corners] at (3.0,-0.95) {Consistent: $h(n)\le c+h(n')$};
\end{scope}
\end{tikzpicture}
\caption{Top: admissibility --- $h$ must lie at or below $h^{*}$. Bottom: consistency as a triangle inequality on every edge.}
\label{fig:admissible}
\end{figure}

\noindent\textbf{Optimality sketch.} With admissible $h$, any suboptimal goal $G_2$ has $f(G_2)=g(G_2)>g^{*}$. An optimal goal $G^{*}$ on the frontier satisfies $f(G^{*})=g^{*}\le f(G_2)$, so A$^{*}$ pops $G^{*}$ first. Dominance matters: if $h_2\ge h_1$ everywhere, A$^{*}$ with $h_2$ expands a subset of the nodes of $h_1$ (up to tie-breaking).

\section{IDA\texorpdfstring{$^{*}$}{*}: Memory-Bounded Optimal Search}
A$^{*}$ keeps the entire frontier ($O(b^{d})$ memory, $b$ branching, $d$ depth). \emph{Iterative-Deepening A$^{*}$} (IDA$^{*}$) runs depth-first search limited by an $f$-threshold, increasing it to the minimum $f$ that exceeded it:
\[
  \text{threshold}_{k+1}=\min\{f(n)\mid f(n)>\text{threshold}_{k}\}.
\]
IDA$^{*}$ is optimal with admissible $h$, uses $O(d)$ memory, but re-expands nodes across iterations --- ideal when memory is scarce and $h$ is strong. In practice, IDA$^{*}$ shines on puzzles (15-puzzle) where $b$ is small.

\begin{table}[h]\centering\small
\begin{tabular}{lccc}
\toprule
Algorithm & $f(n)$ & Optimal? & Memory\\
\midrule
Greedy best-first & $h(n)$ & No & $O(b^{m})$\\
A$^{*}$ (consistent $h$) & $g+h$ & Yes & $O(b^{d})$\\
IDA$^{*}$ (admissible $h$) & $g+h$ threshold & Yes & $O(d)$\\
\bottomrule
\end{tabular}
\caption{Informed search at a glance ($b$ branching, $d$ depth, $m$ max depth).}
\end{table}

\begin{remark}[Why heuristics matter]
A better $h$ can reduce expansions from exponential to linear. On the 8-puzzle, misplaced-tiles expands $\sim$10$\times$ more nodes than Manhattan at depth $20$; on the 15-puzzle the gap is $>100\times$. The cost of computing $h$ must stay small --- an $O(1)$ Manhattan lookup beats a perfect $h^{*}$ that costs exponential time.
\end{remark}

\section{Implementing A\texorpdfstring{$^{*}$}{*}}

\begin{lstlisting}[language=Python]
import heapq

def astar(start, is_goal, successors, heuristic):
    """A* graph search. successors(s) -> [(s', cost)]. Returns (path, cost)."""
    open_heap = [(heuristic(start), 0, start, [start])]  # (f, g, node, path)
    closed = {}  # node -> best g seen
    while open_heap:
        f, g, n, path = heapq.heappop(open_heap)
        if n in closed and closed[n] <= g:
            continue
        closed[n] = g
        if is_goal(n):
            return path, g
        for nxt, c in successors(n):
            ng = g + c
            nf = ng + heuristic(nxt)
            if nxt not in closed or ng < closed.get(nxt, float('inf')):
                heapq.heappush(open_heap, (nf, ng, nxt, path + [nxt]))
    return None, float('inf')

# Example: 4-neighbour grid, cost 1, Manhattan heuristic
def manhattan(a, b):
    return abs(a[0]-b[0]) + abs(a[1]-b[1])
\end{lstlisting}

\begin{lstlisting}[language=Python]
# Heuristic design: when is a heuristic admissible?
# 8-puzzle: tiles 1..8, blank=0, goal has blank bottom-right.

def misplaced_tiles(state, goal):
    """Counts tiles out of place. Admissible: each move fixes at most 1 tile."""
    return sum(1 for a, b in zip(state, goal) if a != 0 and a != b)

def manhattan_8puzzle(state, goal_pos):
    """Sum of Manhattan distances per tile. Dominates misplaced; still admissible."""
    dist = 0
    for idx, tile in enumerate(state):
        if tile == 0:
            continue
        r, c = divmod(idx, 3)
        gr, gc = goal_pos[tile]
        dist += abs(r - gr) + abs(c - gc)
    return dist

# Consistency check for an edge n -> n' with cost c:
#   assert heuristic(n) <= c + heuristic(n_prime)
# Straight-line distance on a map passes this (triangle inequality).

# Admissible but inconsistent example (graph search must reopen):
#   h(A)=4, h(B)=1, edge A->B cost 1 gives 4 <= 1+1? No -> inconsistent.
\end{lstlisting}

\section{Choosing and Combining Heuristics}
A heuristic comes from a \emph{relaxed problem}: drop constraints and solve cheaply. Straight-line distance relaxes roads to Euclidean plane; Manhattan relaxes diagonal moves; misplaced-tiles relaxes that tiles interfere. If $h_1,h_2$ are admissible, so is $h=\max\{h_1,h_2\}$ and $h$ dominates both (expands no more nodes, up to ties). Pattern databases precompute exact costs for subproblems (e.g., 7 tiles of the 15-puzzle) and remain admissible by partitioning moves.
The \emph{effective branching factor} $b^{*}$ solves $N+1=1+b^{*}+(b^{*})^{2}+\dots+(b^{*})^{d}$ for $N$ expanded nodes and depth $d$; better $h$ drives $b^{*}$ toward $1$.

\section{Local Search and Metaheuristics}
Sometimes the goal is a good state, not a path: placing 8 queens with none attacking,
scheduling exams into few slots, or tuning weights. Storing a whole frontier is wasteful;
keep one state and improve it. \emph{Hill climbing} moves to the best-valued neighbour
and stops when none is better. It is fast but sticks at local optima, plateaus, and ridges.
Random restarts (many climbs, keep the best) or allowing sideways moves mitigates this.
\emph{Simulated annealing} escapes local optima by accepting a worse move with probability
$e^{\Delta E/T}$, where $\Delta E < 0$ measures the loss and $T$ is a temperature cooled
over time. High $T$ explores; as $T \to 0$ the method becomes hill climbing. A slow cooling
schedule converges to a global optimum with probability approaching one.
\emph{Genetic algorithms} keep a population of states. Fitter states reproduce: crossover
splices two parents and mutation perturbs offspring, so useful building blocks combine.
Table~\ref{tab:local-search} contrasts the three.
\begin{table}[h]\centering\small
\begin{tabular}{lcc}
\toprule Method & State kept & Escapes local optima \\\midrule
Hill climbing & one & restarts only \\
Simulated annealing & one & worse moves w.p.\ $e^{\Delta E/T}$ \\
Genetic algorithm & population & crossover + mutation \\\bottomrule
\end{tabular}
\caption{Local search at a glance: one state or many, and how each escapes local optima.}
\label{tab:local-search}
\end{table}
\section{Key Takeaways}
Greedy ($f=h$) is fast but not optimal; A$^{*}$ ($f=g+h$) is optimal with admissible $h$ and never reopens with consistent $h$; IDA$^{*}$ trades time for $O(d)$ memory via $f$-thresholds. Heuristic quality dominates performance: a dominated heuristic can expand exponentially more nodes. Design $h$ by relaxing the problem (e.g., ignore obstacles for straight-line distance) and prefer the dominating admissible heuristic (e.g., $\max\{h_1,h_2\}$). When memory is tight, IDA$^{*}$ with a strong heuristic beats A$^{*}$; when $h$ is weak, A$^{*}$'s best-first order saves re-expansions.

\section{Exercises}
\begin{exercise}Prove that consistency implies admissibility. Give a 3-node graph where $h$ is admissible but not consistent.\end{exercise}
\begin{exercise}Run A$^{*}$ on Fig.~\ref{fig:astar-tree} listing the order nodes are popped from the frontier. What would greedy best-first ($f=h$) expand first?\end{exercise}
\begin{exercise}Show that if $h_1$ and $h_2$ are admissible then $h(n)=\max\{h_1(n),h_2(n)\}$ is admissible and dominates both. When is $\max$ still consistent?\end{exercise}
\begin{exercise}For the 8-puzzle, compare misplaced-tiles vs.\ Manhattan distance on a state where tiles 4 and 5 are swapped. Compute both values and explain which prunes more.\end{exercise}
\begin{exercise}IDA$^{*}$ starts with threshold $=h(S)$. Trace two iterations on a line $S\to A\to G$ with costs $1,1$ and $h(S)=2,h(A)=2,h(G)=0$. What is the second threshold? Why does IDA$^{*}$ re-expand $S$?\end{exercise}
